
% ════════════════════════════════════════
% 模型的改进与推广（08.模型改进推广.tex）
% ════════════════════════════════════════
\section{模型的改进与推广}

\subsection{模型的改进}

本文模型的不足主要体现在三个方面。在方法层面，可将光伏预报误差建模为随机变量，引入两阶段随机规划或鲁棒优化，在目标函数中显式权衡期望购电费用与违约风险，使调整策略更具前瞻性。在参数层面，可将储能充放电效率进一步细分为充电效率与放电效率、考虑电池循环寿命损耗与自放电率，并引入储电量分段折旧成本，使经济模型更贴近实际设备。在数据层面，可结合历史气象数据与数值天气预报，对光伏预报误差的分布进行估计，据此优化各预报时刻的储能安全裕度，进一步降低紧急购电费用。

\subsection{模型的推广}

本文建立的微网购电与储能调控模型具有较强的通用性，可直接推广到更广泛的工程与政策场景。在工程领域，该模型可用于工业园区、商业综合体与电动汽车充电站的储能容量配置与日前调度，也可嵌入虚拟电厂的能量管理系统，实现多主体分布式资源的协同优化；在政策领域，该模型可用于评估分时电价、尖峰电价等价格机制的引导效果，为需求响应与储能补贴政策的制定提供量化依据；在行业领域，其滚动调整与惩罚机制思想可迁移到电力现货市场交易、天然气购销合同与物流库存管理等存在预测更新与违约成本的场景，具有广泛的适用价值。
